\section{Actual macrostep symmetries and the quadratic tensor}
\label{msy-section}

This section concerns the actual local composite sector and its
common-domain analytic continuation. It does not assert equality of
different gauge sectors or spatial symmetry of an arbitrary measurable
completion outside the analytic domain. The input total density must
have the spatial symmetries under consideration.

\subsection{Symmetries of the exact local integral}
\label{msy-integral-section}

Fix the odd primitive ratio \(\ell\), a finite power \(R=\ell^s\), and
the centered nested blocks of the composite construction. Use compatible
\emph{cubic} tori, with distinct lifted interfaces as in the geometric
proof. The sufficiently large auxiliary grouping \(J\) can be chosen
as a power of \(\ell\), preserving a cofinal family of source tori.
On rectangular tori only permutations preserving the side lengths
are available.

The selected source mean is assumed to obey its permutation axiom
as well as the endpoint covariance already used in the construction.
The all-coordinate-order contour family and this axiom are the
properties stated in Ba{\l}aban, CMP 109 (1987), pp.~252--254,
especially (0.7). The implicit mean (0.10) is an admissible example;
it is not the definition of every permitted mean.

\begin{proposition}[Exact spatial symmetry of the local sector]
\label{msy-spatial-proposition}
Suppose the complete input density is invariant under the relevant
spatial transformations on the common small domain. The composite
constraint, independent direct-\(R\) internal contour gauge, physical
centered branch and actual retained link/color-ball integral commute
with
\begin{equation}
 x\longmapsto Sx+Rn,\qquad
 n\in\mathbb Z^d,\quad
 S\text{ a signed coordinate permutation}.
 \label{msy-space-group}
\end{equation}
After the output lattice is rescaled by \(R\), these give all its unit
translations and the full hypercubic point group. The same symmetries
hold for the total generated correction on the common zero-free
holomorphic domain, with its compatible logarithm branch.
\end{proposition}

\begin{proof}
Every transformation in \eqref{msy-space-group} preserves each level
of the centered hierarchy. It permutes the coordinate-order contour
family. The mean's permutation axiom therefore makes the contour
maps equivariant. Orientation reversal of an ordinary coarse bond
also respects the averaging formula. If that formula is written as
\[
 F=\exp\!\left\{N^{-1}\sum_x\log(U_xT^{-1})\right\}T,
\]
the reversed strings and reference are \(U_x^{-1},T^{-1}\), and
\begin{equation}
 \log(U_x^{-1}T)
     =-\operatorname{Ad}_{T^{-1}}\log(U_xT^{-1}),
 \qquad F_{\rm reverse}=F^{-1}.
 \label{msy-average-reversal}
\end{equation}
The primitive maps, their composition and the separate direct-\(R\)
contour map consequently commute with the transformation. The
centerline interface pivots are preserved as unoriented bonds.

The relative-coordinate calculation is essential. For
\(U_b=e^{iB_b}V_b\), a reversed bond satisfies
\begin{equation}
 U_b^{-1}
   =e^{-i\operatorname{Ad}_{V_b^{-1}}B_b}V_b^{-1}.
 \label{msy-relative-reversal}
\end{equation}
At fixed real \(V\), this is an orthogonal linear color map.
Together with the bond permutation it induces an orthogonal map
\(D_{S,V}\) on the retained coordinates. Uniqueness of the small
pivot solution gives
\begin{equation}
 \Phi_{SV}(D_{S,V}z)
       ={\cal R}_{S,V}\Phi_V(z),
 \label{msy-chart}
\end{equation}
where \({\cal R}_{S,V}\) is the full relative-bond transformation.
Thus the original window
\(\prod_{b\ {\rm retained}}1_{\{|z_b|\le gT\}}\)
is preserved with absolute Jacobian one. No complex change of a
sharp indicator is used.

The transformed physical branch is feasible, has all internal
contours equal to the identity, and has the same constrained
classical value and stationarity. Common-domain uniqueness yields
\(V_*(SW)=SV_*(W)\). The homogeneous classical action has the
required value symmetry by applying the same transformations
to its defining ordinary-constraint minimization; the fine Wilson
action and small global norm class are preserved.

Haar factors are invariant under inverses and conjugation. A pivot
derivative transforms by the relative coarse and pivot color maps.
On a reversed bond both have determinant
\((-1)^{\dim\mathfrak g}\); their adjoint factors have determinant
one. The signs cancel in the ratio. Permutations merely relabel
the factors. This accounts for both the background delta
determinant and the normalized straightening determinant.

The chart action plus absolute gauge penalty satisfies
\(T_{SV}(D_{S,V}z)=T_V(z)\). The map \(D_{S,V}\) is independent
of \(z\), so constant, linear and quadratic Taylor terms transform
separately, also in the centered independent-background extension.
The Hessian transforms by orthogonal congruence, preserving the
normalized Gaussian and its determinant. The fixed flat determinant
is the same scalar denominator. The retained change of variables
therefore proves equality of the real local integrals.

The complete inherited insertion agrees with its pre-gauge
representation on the original common domain. Its assumed
spatial invariance and holomorphic uniqueness extend the total
identity to the common-domain continuation. No symmetry of its
individual activity coefficients is required. Finally the proved
zero-free cluster domain and common branch give equality of
logarithms. Holomorphy alone would not imply this last step.
\end{proof}

Auxiliary \(J\)-grouping reorganizes the exact retained product
window and finite integral; it does not alter their scalar total.
Its fixed carrier grid and deterministic trees may break symmetry
at coefficient level. The proposition concerns their sum.
Additional genuine block-norm or asymmetric cutoff restrictions
would require a new check. A color-coordinate cube is not a
substitute for the invariant linkwise color balls.

There is an explicit obstruction to current-lattice unit
translations for the fixed sector. A pure gauge with a small
nonzero phase \(t\) at one block root and zero elsewhere has
absolute contour penalty
\[
 (R^d-1)q(e^{it}),\qquad
 q(U)=1-\tfrac12\tau(U+U^{-1}).
\]
Moving that phase one current unit to a nonroot gives only
\(q(e^{it})\). This does not conflict with output unit translations,
whose current lift is \(R\).

The measurable completion using an ordered reference contour and
fallback is gauge equivariant, but its global spatial equivariance
has not been proved. The proposition does not concern that
completion or the uncontrolled complement. If needed, averaging
transformed sector--complement decompositions preserves their
exact sum when the blocking map and input are equivariant.
Averaging conjugated asymmetric blocking kernels instead defines
a new Markov blocking kernel; it is not equality with the original
deterministic map. No such change is made here.

\paragraph{Charge conjugation.}
For \(SU(2)\), \(U\mapsto\overline U\) is conjugation by
\(i\sigma_2\), hence is already a constant gauge transformation.
It is exact on the same domains. For \(SU(N)\), the holomorphic
automorphism is \(U\mapsto(U^{-1})^T\), with orthogonal real Lie
map \(A\mapsto-A^T\). The implicit mean
\(\sum_j\log(U_jP^{-1})=0\) commutes with this automorphism
by local uniqueness. The proof then applies if the complete
input is also invariant. For a different unspecified allowed
mean, this outer-automorphism covariance is an extra property
to verify; reality alone is not a proof.

\subsection{Uniform Hessian decay and a kernel limit}
\label{msy-kernel-section}

Let the \emph{total} output correction have connected activities
on output \(J\)-cells:
\begin{equation}
 F(w)=\sum_{H\ {\rm connected}}F_H(w_H),\qquad
 \sup_I\sum_{H\ni I}e^{a_0|H|}\|F_H\|_b\le\epsilon.
 \label{msy-activity-norm}
\end{equation}
Here the supremum is on a scalar-coordinate polydisc of radius
\(b\); a color-ball domain may be replaced by a contained scalar
polydisc. Write \(D_0=d(\dim\mathfrak g)J^d\) for an upper bound
on the scalar coordinate count in a cell. All parameters are
uniform in volume and preceding homogeneous depth; \(J\) is fixed.

For bond/color coordinates \(i,j\), let \(\xi_i,\xi_j\) be their
bond midpoints and \(K_{ij}=\partial_i\partial_jF(0)\).
Cauchy gives
\[
 |\partial_i\partial_jF_H(0)|\le2b^{-2}\|F_H\|_b.
\]
If connected \(H\) contains both cells, their midpoint distance
is at most \(J(|H|+d)\), using shortest distance on a torus.
Sum the at most \(D_0|H|\) second coordinates before the
anchored support sum. With \(\mu_0=a_0/(2J)\), this proves
\begin{equation}
 \sup_i\sum_j e^{\mu_0|\xi_i-\xi_j|_1}|K_{ij}|
 \le {\cal M}:=
 \frac{2D_0e^{a_0d/2}\epsilon}{b^2}
                  \sup_{n\ge1}n e^{-a_0n/2}.
 \label{msy-exponential-row}
\end{equation}
No symmetry of the activity decomposition was used.

For the real quadratic classification assume \(F\) is real on real
small fields, as for the physical positive density and its real
logarithm. Complex inputs have the same algebraic classification
with a possibly complex \(\beta\).
For a compact semisimple group, constant gauge invariance
gives \(DF(0)=0\). Differentiate local gauge invariance once
in a site gauge and once in a field direction. The derivative
of its nonlinear generator is multiplied by \(DF(0)\), so
the exact position-space Ward identity is
\begin{equation}
 K D=0,
 \label{msy-Ward-realspace}
\end{equation}
with \(D\) the output unit gradient.

Take an exhausting sequence of compatible cubic tori and
assume the uniform bound \eqref{msy-exponential-row}.
A diagonal subsequence gives a kernel limit at every fixed
displacement. Fatou preserves exponential summability.
Translation covariance, signed permutations, Hessian symmetry
and the finite-difference identity \eqref{msy-Ward-realspace}
pass at fixed displacements. This supplies an infinite-lattice
kernel, not a probability measure, uniqueness statement or
continuum limit.

This limiting step is substantive. Finite periodic Ward
identities hold only at their allowed momenta and do not
by themselves imply a vanishing zero-momentum mass
coefficient. Small gauge-invariant winding terms are possible.
No interpolation of finite discrete momenta is used below.

\subsection{Midpoint Fourier classification with a remainder bound}
\label{msy-classification-section}

Constant adjoint invariance makes the color structure scalar
on each simple ideal. Cross terms between distinct simple
ideals vanish. Work with one such spatial kernel
\(k_{\mu\nu}(r)\). Define
\begin{equation}
 \widetilde K_{\mu\nu}(p)=
 \sum_r k_{\mu\nu}(r)
       e^{ip\cdot(r+(e_\nu-e_\mu)/2)},\qquad
 \widehat p_\mu=2\sin(p_\mu/2).
 \label{msy-midpoint}
\end{equation}
The exponential bound makes the symbol analytic near zero.
The limiting position-space Ward identity yields
\begin{equation}
 \widetilde K(p)\widehat p=0
       \quad\text{for every real }p.
 \label{msy-Ward-symbol}
\end{equation}
At \(p=te_\mu\ne0\), this kills \(e_\mu\); continuity gives
\(\widetilde K(0)=0\). Full inversion gives
\(\widetilde K(-p)=\widetilde K(p)\). Reality and Hessian
symmetry make the symbol real symmetric for real \(p\).
Midpoint coordinates give the phase-free signed-permutation
law and eliminate odd Taylor degrees.

\begin{theorem}[Quadratic Maxwell tensor]
\label{msy-Maxwell-theorem}
Under the preceding total symmetries, Ward identity and
uniform kernel hypotheses, each simple color factor has
a real coefficient \(\beta\) such that
\begin{equation}
 \widetilde K(p)
 =\beta\bigl(|p|^2I-pp^T\bigr)+O(|p|^4)
 =\beta\bigl(|\widehat p|^2I-
                      \widehat p\widehat p^T\bigr)+O(|p|^4).
 \label{msy-Maxwell-expansion}
\end{equation}
The estimate concerns the subsequential infinite kernel;
it is not an assertion about a continuous interpolation
of a finite-volume symbol.
\end{theorem}

\begin{proof}
Signed reflections force the degree-two part to have only
squares on its diagonal and \(p_\mu p_\nu\) off diagonal.
Permutations identify their coefficients:
\[
 (K_2)_{\mu\mu}=A p_\mu^2+
                  B\sum_{\nu\ne\mu}p_\nu^2,\qquad
 (K_2)_{\mu\nu}=C p_\mu p_\nu\quad(\mu\ne\nu).
\]
The cubic part of \eqref{msy-Ward-symbol} is \(K_2(p)p=0\).
Its polynomial coefficients give \(A=0\) and \(C=-B\).
Put \(\beta=B\). Analyticity and parity give the fourth-order
remainder, and \(\widehat p=p+O(|p|^3)\) gives its second form.
\end{proof}

For explicit constants set
\begin{equation}
 M_n=\max_\mu\sum_{\nu,r}
 |r+(e_\nu-e_\mu)/2|_1^n\,|k_{\mu\nu}(r)|.
 \label{msy-moments}
\end{equation}
For \(n>0\), \eqref{msy-exponential-row} gives
\[
 M_n\le{\cal M}\left(\frac{n}{e\mu_0}\right)^n,
 \qquad |\beta|\le M_2/2.
\]
The even cosine Taylor remainder yields, for real \(p\),
\[
 \|\widetilde K(p)-K_2(p)\|_{\rm op}
               \le M_4\|p\|_\infty^4/24.
\]
Using \(|2\sin(t/2)-t|\le|t|^3/24\) in every row,
for \(\|p\|_\infty\le1\) we obtain
\begin{equation}
 \left\|\widetilde K(p)
       -\beta\bigl(|\widehat p|^2I-
                        \widehat p\widehat p^T\bigr)\right\|_{\rm op}
 \le\left(\frac{M_4}{24}+\frac d6|\beta|\right)
                              \|p\|_\infty^4 .
 \label{msy-quantitative-remainder}
\end{equation}
Real symmetry turns the absolute row estimate into the
operator estimate. Without coordinate permutations there
need not be a single \(\beta\): unequal plaquette-orientation
coefficients already give a local gauge-invariant action
with translations and axis reflections but several
quadratic coefficients.

\subsection{Scope for a subsequent extraction}
\label{msy-scope-section}

The preceding argument establishes actual output-lattice
symmetry of the local \(SU(2)\) integral and the resulting
quadratic classification. It determines neither the sign
of \(\beta\) nor asymptotic freedom. The uniform flat
classical normalization is a separate, already available
input: the explicit alias formula in Notebook07 gives
\[
 q_{M,\mathrm{mid}}(p)=|p|^2I-pp^T+O(|p|^4)
\]
uniformly in \(M\). Together with the homogeneous minimizer's
identification of \(q_M\) as \(D^2S_M(0)\), this supplies
the classical quadratic normalization in that scope.
Ellipticity alone would not have supplied it.

Subtracting \(\beta\) times the normalized standard Wilson
plaquette functional is a gauge-invariant scalar subtraction
whose residual quadratic kernel satisfies
\eqref{msy-quantitative-remainder}. A localized nonlinear
extraction, the higher-field Ward completion, finite-volume
winding control and stable relevant/marginal flow remain
additional steps. Coefficientwise high-field-jet
contraction by itself does not supply them.
